% chapter09.tex - 等效原理与广义协变性
\chapter{等效原理与广义协变性 (Equivalence Principle and General Covariance)}
\label{chap:equivalence}

等效原理是广义相对论最核心的物理基石，广义协变性原理则为理论提供了数学框架。本章系统阐述这两大原理的物理内涵、实验检验和数学表述，并由此自然导出弯曲时空中的物理定律构造方法——最小耦合原理。

\section{引力与加速度的等价性}

\subsection{升降机思想实验}

爱因斯坦在1907年提出了著名的升降机思想实验，深刻揭示了引力与加速度之间的内在联系。考虑两个完全封闭的升降机：

\begin{enumerate}
    \item \textbf{升降机 A}：静止在地球表面，感受到向下的均匀引力场 $\vb*{g} = -g \vu*{z}$。
    \item \textbf{升降机 B}：在无引力的空无一物的空间中，以恒定加速度 $\vb*{a} = +g \vu*{z}$ 向上加速。
\end{enumerate}

\begin{figure}[htbp]
\centering
\begin{tikzpicture}[scale=1.1]
    % --- 升降机 A：地球引力 ---
    \draw[thick] (-2.5, -1.5) rectangle (0.5, 2.0);
    \node at (-1.0, 1.7) {升降机 A};
    % 内部的人
    \fill (-1.0, 0.3) circle (0.15);
    \draw (-1.0, 0.3) -- (-1.0, -0.3);
    \draw (-1.0, -0.3) -- (-1.2, -0.8);
    \draw (-1.0, -0.3) -- (-0.8, -0.8);
    % 掉落的苹果
    \fill[red] (-0.5, 1.2) circle (0.08);
    \draw[red, ->, thick] (-0.5, 1.2) -- (-0.5, 0.5);
    \node[red, right] at (-0.5, 0.8) {苹果下落};
    % 地球
    \shade[ball color=brown!40] (-4.0, -3.0) circle (0.8);
    \draw[->, thick] (-1.0, -1.5) -- (-1.0, -1.9);
    \node at (-1.0, -2.15) {$\vb*{g}$ 向下};
    % 标签
    \node at (-1.0, -2.6) {引力场中静止};

    % --- 升降机 B：加速参考系 ---
    \begin{scope}[xshift=4.5cm]
    \draw[thick] (-2.5, -1.5) rectangle (0.5, 2.0);
    \node at (-1.0, 1.7) {升降机 B};
    % 内部的人（漂浮）
    \fill (-1.0, 0.3) circle (0.15);
    \draw (-1.0, 0.3) -- (-1.0, -0.3);
    \draw (-1.0, -0.3) -- (-1.2, -0.8);
    \draw (-1.0, -0.3) -- (-0.8, -0.8);
    % 苹果相对人"下落"
    \fill[red] (-0.5, 1.2) circle (0.08);
    \draw[red, ->, thick] (-0.5, 1.2) -- (-0.5, 0.5);
    \node[red, right] at (-0.5, 0.8) {苹果"下落"};
    % 加速箭头
    \draw[->, thick] (-1.0, -1.5) -- (-1.0, -1.9);
    \node at (-1.0, -2.15) {$\vb*{a}$ 向上};
    \node[blue] at (-1.0, -2.7) {加速上升};
    % 虚线表示无引力空间
    \draw[blue, dashed] (-0.2, -3.3) -- (3.2, -3.3);
    \node[blue] at (1.5, -3.5) {无引力真空};
    \end{scope}

    % 等号
    \node at (2.0, 0.25) {\Huge $\Longleftrightarrow$};
    \node at (2.0, -0.4) {不可分辨};

\end{tikzpicture}
\caption{升降机思想实验。升降机A静止在引力场中，苹果因引力下落；升降机B在无引力空间中加速上升，苹果因惯性而相对于加速参考系"下落"。封闭在升降机内的观察者无法通过任何局部力学实验区分这两种情况。}
\label{fig:elevator}
\end{figure}

\textbf{关键结论：}封闭在升降机内部的观察者无法通过任何局部实验来判断自己处于引力场中还是处于加速参考系中。这一观察可以精确表述为：

\begin{definition}[弱等效原理 (Weak Equivalence Principle, WEP)]
在引力场中，所有检验物体的运动轨迹仅取决于其初始位置和速度，而与物体的质量和内部组成无关。等价的数学表述为：引力质量 $m_G$ 与惯性质量 $m_I$ 相等，
\begin{equation}
    m_I = m_G .
\end{equation}
\label{def:WEP}
\end{definition}

\subsection{E\"{o}tv\"{o}s 实验与现代检验}

弱等效原理的实验基础可以追溯到伽利略的比萨斜塔思想实验和牛顿的单摆实验，但系统性的精密检验始于E\"{o}tv\"{o}s。

\begin{example}[E\"{o}tv\"{o}s 扭摆实验，1890--1922]
E\"{o}tv\"{o}s 扭摆实验的原理如下：将两个不同材料（如木和铂）的物体悬挂在扭摆的两端。若 $m_G/m_I$ 因材料而异，则地球自转产生的离心力与引力的合力方向将随材料变化，从而在扭摆上产生净扭矩。设地球引力加速度为 $\vb*{g}$，自转离心加速度为 $\vb*{a}_c = \omega^2 R_\oplus \sin\theta \, \vu*{e}_\perp$，则悬挂方向由下式决定：
\begin{equation}
    \tan\alpha = \frac{m_I}{m_G} \frac{a_c}{g}.
\end{equation}
若 $m_I/m_G$ 依赖于材料成分，则 $\alpha$ 随材料变化，扭摆会转动。E\"{o}tv\"{o}s 得到
\begin{equation}
    \left|\frac{m_G}{m_I} - 1\right| < 5 \times 10^{-9}.
\end{equation}
\end{example}

现代实验已将这一精度推向极致：

\begin{itemize}
    \item \textbf{E\"{o}t-Wash 实验}（华盛顿大学，1990--2008）：利用改进的扭摆技术，达到
    \begin{equation}
        \left|\frac{m_G}{m_I} - 1\right| < 10^{-13}.
    \end{equation}
    \item \textbf{MICROSCOPE 卫星}（法国，2016--2022）：在太空微重力环境中进行，精度达到
    \begin{equation}
        \eta = \frac{2(m_G/m_I)_A - 2(m_G/m_I)_B}{(m_G/m_I)_A + (m_G/m_I)_B} < 2.7 \times 10^{-15}.
    \end{equation}
    \item \textbf{月球激光测距}（LLR）：通过精确测量月球轨道，检验了地球和月球的 $m_G/m_I$ 等价性，精度约 $10^{-13}$。
\end{itemize}

这些实验使弱等效原理成为物理学中最精确检验的原理之一。

\subsection{等效原理的三层表述}

等效原理并非单一陈述，而是从弱到强的三层嵌套表述：

\begin{definition}[弱等效原理 (WEP)]
引力质量等于惯性质量，检验粒子在引力场中的运动与组成无关。
\end{definition}

\begin{definition}[爱因斯坦等效原理 (Einstein Equivalence Principle, EEP)]
WEP 加上以下两个条件：
\begin{enumerate}
    \item \textbf{局部 Lorentz 不变性 (LLI)}：在任何局部自由下落参考系中，非引力的物理实验结果与参考系的速度（方向和大小）无关。
    \item \textbf{局部位置不变性 (LPI)}：在任何局部自由下落参考系中，非引力的物理实验结果与实验在时空中的位置和时间无关。
\end{enumerate}
简言之，在局部自由下落参考系中，狭义相对论（以及整个非引力的标准模型物理）成立。
\end{definition}

\begin{definition}[强等效原理 (Strong Equivalence Principle, SEP)]
EEP 再加一个条件：引力自身能量的 $m_G/m_I$ 比值也等于1。换言之，\textbf{所有物理定律}——包括引力定律本身——在局部自由下落参考系中取其在狭义相对论中的形式。这意味着引力的自能量（self-energy）不产生异常的惯性或引力效应。
\end{definition}

\begin{note}[SEP 与 WEP/EEP 的区别]
广义相对论满足 EEP，但是否满足 SEP 取决于具体情况。在纯广义相对论中（无标量场等），SEP 成立。但在许多修正引力理论（如标量-张量理论）中，引力自能的 $m_G/m_I$ 可能偏离1，导致 SEP 破缺。因此 SEP 的检验可以区分广义相对论与其他引力理论。
\end{note}

\begin{table}[htbp]
\centering
\caption{等效原理的三层表述及其检验}
\begin{tabular}{p{2.5cm} p{5cm} p{5cm}}
\toprule
\textbf{原理} & \textbf{核心内容} & \textbf{检验实验} \\
\midrule
WEP & $m_I=m_G$，运动轨迹与组成无关 & E\"{o}tv\"{o}s, E\"{o}t-Wash, MICROSCOPE \\
EEP & 局部物理定律与SR一致 & 引力红移 (Pound-Rebka)、Hughes-Drever实验 \\
SEP & 引力自能量也满足 $m_I=m_G$ & 月球激光测距（Nordtvedt效应）、脉冲星双星 \\
\bottomrule
\end{tabular}
\label{tab:ep_levels}
\end{table}

\section{引力时间膨胀的推导}

\subsection{从加速参考系的多普勒频移到引力红移}

等效原理最深刻的推论之一是引力场中的时间膨胀效应。我们可以通过加速参考系中的多普勒频移来推导。

\textbf{步骤一：加速参考系中的频移。} 考虑一个在无引力空间中沿 $z$ 轴以恒定加速度 $a$ 加速的火箭。在火箭的瞬时静止参考系中，等效原理告诉我们它与均匀引力场中的静止参考系不可区分。设在火箭底部（$z=0$）有一个光源，顶部（$z=h$）有一个探测器。

在 $t=0$ 时刻，底部光源发出一个频率为 $\nu_0$ 的光信号。此时火箭的瞬时速度为0。光信号以速度 $c$ 向上传播，到达顶部的时间为 $\Delta t \approx h/c$。在这段时间内，火箭获得的速度增量为
\begin{equation}
    \Delta v = a \Delta t = \frac{a h}{c}.
\end{equation}
由于多普勒效应，顶部探测器接收到的频率为（对 $v \ll c$ 的非相对论情况）：
\begin{equation}
    \nu_{\text{rec}} = \nu_0 \left(1 - \frac{\Delta v}{c}\right) = \nu_0 \left(1 - \frac{ah}{c^2}\right).
\end{equation}

\textbf{步骤二：等效原理翻译为引力红移。} 根据等效原理，加速参考系中的物理等效于均匀引力场中的物理。将加速度 $a$ 替换为引力加速度 $g$（即单位质量所受引力 $g = GM/r^2$），得到：
\begin{equation}
    \frac{\nu_{\text{rec}} - \nu_0}{\nu_0} = -\frac{g h}{c^2}.
\end{equation}
这意味着在引力场中，从低引力势向高引力势传播的光会发生红移。

\begin{theorem}[引力红移公式——弱场近似]
在弱静态引力场中，设 $\Phi(\vb*{x})$ 为牛顿引力势（满足 $\nabla^2 \Phi = 4\pi G \rho$），频率变化为
\begin{equation}
    \boxed{\frac{\Delta\nu}{\nu} = \frac{\Phi_{\text{em}} - \Phi_{\text{obs}}}{c^2}}.
\end{equation}
等价地，坐标时与固有时的关系为
\begin{equation}
    \dd\tau = \sqrt{1 + \frac{2\Phi}{c^2}} \, \dd t \approx \left(1 + \frac{\Phi}{c^2}\right) \dd t.
\end{equation}
\end{theorem}

\begin{proof}[加速参考系推导的严格版本]
设火箭底部光源发射光信号的时刻为 $t_1$（火箭瞬时静止），顶部探测器接收到信号的时刻为 $t_2$。在 Rindler 坐标系中，火箭底部的固有时 $d\tau_{\text{bottom}}$ 与顶部的固有时 $d\tau_{\text{top}}$ 的关系为
\begin{equation}
    d\tau_{\text{top}} = \left(1 + \frac{ah}{c^2}\right) d\tau_{\text{bottom}}.
\end{equation}
将 $a \to g$ 翻译为引力场，即得 $d\tau_{\text{top}}/d\tau_{\text{bottom}} = 1 + gh/c^2$。对于地球表面，
\begin{equation}
    \frac{gh}{c^2} \approx \frac{9.8 \times 10}{9 \times 10^{16}} \approx 1.1 \times 10^{-15}.
\end{equation}
这是一个极小的效应，但可以被精密测量所检测。
\end{proof}

\subsection{Pound--Rebka 实验 (1960)}

\begin{example}[Pound--Rebka 实验]
1960年，Pound 和 Rebka 在哈佛大学 Jefferson 实验室的塔楼中（高度差 $h=22.5\,\text{m}$）利用 M\"{o}ssbauer 效应测量了引力红移。$^{57}\text{Fe}$ 的 14.4 keV $\gamma$ 射线具有极窄的线宽（$\Delta\nu/\nu \approx 10^{-12}$），使如此微小的频移可被检测。

理论预测的频移为
\begin{equation}
    \frac{\Delta\nu}{\nu} = \frac{gh}{c^2} = \frac{9.8 \times 22.5}{9 \times 10^{16}} \approx 2.46 \times 10^{-15}.
\end{equation}
实验结果与理论预测的偏差在 $10\%$ 以内，后经 Pound--Snider (1964) 改进至 $1\%$。

1976年的 \textbf{Gravity Probe A} 实验将氢钟搭载到 $10^4\,\text{km}$ 高度，以 $7\times 10^{-5}$ 的精度验证了引力红移公式。2018年的 \textbf{GALILEO 5/6 卫星}（因发射故障进入椭圆轨道）提供了 $10^{-4}$ 精度的额外检验。
\end{example}

\section{光线偏折的等效原理推导}

\subsection{等效原理预言光线偏折}

等效原理不仅能推导引力红移，还能预言光线在引力场中的偏折。虽然等效原理给出的偏折角恰好是广义相对论精确值的一半，但这个推导极具教学意义——它清晰地展示了等效原理的适用范围和局限。

\begin{example}[用等效原理推导光线偏折——因子2错误]
考虑一束光从远处掠过太阳表面。在太阳附近，引力场可近似为均匀的。设光子在太阳引力场中沿 $x$ 方向传播，引力沿 $z$ 方向。

\textbf{等效原理推理：} 在自由下落参考系中，光子沿直线传播。从地球（远离太阳的惯性系）看来，太阳附近的时空可以看作一个"加速向下"的区域。光子穿越这一区域的时间约为 $\Delta t \approx 2R_\odot / c$（直径穿越），在此期间光子获得的横向速度增量为
\begin{equation}
    \Delta v_z = g \Delta t = \frac{GM_\odot}{R_\odot^2} \cdot \frac{2R_\odot}{c} = \frac{2GM_\odot}{c R_\odot}.
\end{equation}
因此偏折角为（小角度近似）
\begin{equation}
    \boxed{\Delta\theta_{\text{EP}} = \frac{\Delta v_z}{c} = \frac{2GM_\odot}{c^2 R_\odot}}.
\end{equation}
代入太阳数据：$G = 6.674 \times 10^{-11}\,\text{m}^3\text{kg}^{-1}\text{s}^{-2}$, $M_\odot = 1.989 \times 10^{30}\,\text{kg}$, $R_\odot = 6.957 \times 10^8\,\text{m}$, $c = 2.998 \times 10^8\,\text{m/s}$，
\begin{equation}
    \Delta\theta_{\text{EP}} \approx 0.875'' \;\;(\text{角秒}).
\end{equation}
\end{example}

\begin{note}[因子2的缺失]
广义相对论给出的精确偏折角是 \textbf{两倍}于等效原理的预言：
\begin{equation}
    \boxed{\Delta\theta_{\text{GR}} = \frac{4GM_\odot}{c^2 R_\odot} \approx 1.75''}.
\end{equation}
等效原理漏掉的因子2来源于空间弯曲的贡献。等效原理只考虑了"时间弯曲"（$g_{00}$ 分量）——即引力时间膨胀对光速的影响，而没有计入空间弯曲（$g_{ij}$ 分量）对光传播路径的额外偏折。在广义相对论中，时间弯曲和空间弯曲各贡献 $\Delta\theta/2 = 0.875''$。这一事实说明：\textbf{等效原理在局部成立，但光线偏折是一个非局部的、整体的效应，需要完整的时空弯曲描述。}
\end{note}

\subsection{Eddington 日食观测 (1919) 与后续检验}

\begin{example}[Eddington 日食观测]
1919年5月29日，Arthur Eddington 率领两支远征队分别前往巴西的 Sobral 和西非的 Pr\'{i}ncipe 岛观测日全食。通过比较日食期间和夜晚拍摄的星场照片，测量星光掠过太阳边缘时的偏折角。结果：
\begin{equation}
    \Delta\theta_{\text{Sobral}} = 1.98'' \pm 0.16'', \quad
    \Delta\theta_{\text{Pr\'{i}ncipe}} = 1.61'' \pm 0.40''.
\end{equation}
虽然误差较大，但结果明确支持广义相对论（$1.75''$）而非牛顿预言（$0.875''$，等效原理值）。

\textbf{现代射电干涉测量：} 甚长基线干涉测量（VLBI）已将光线偏折的测量精度提升至 $10^{-4}$ 量级，完全确认了广义相对论的 $1.75''$ 预言。
\end{example}

\section{广义协变性原理}

\subsection{数学表述}

广义协变性原理是广义相对论的基本对称性原则：

\begin{definition}[广义协变性原理 (Principle of General Covariance)]
物理定律必须用\textbf{张量方程}表达，使其在任意坐标变换下保持形式不变。具体而言，若物理定律在坐标系 $\{x^\mu\}$ 中写为
\begin{equation}
    T^{\mu_1\cdots\mu_k}_{\nu_1\cdots\nu_\ell} = 0,
\end{equation}
则在任意坐标变换 $x^\mu \to x'^\mu$ 下，它自动变为
\begin{equation}
    T'^{\mu_1\cdots\mu_k}_{\nu_1\cdots\nu_\ell} = 0.
\end{equation}
张量变换保证了这一性质。
\end{definition}

在数学上，广义协变性等价于物理定律在微分同胚群 $\mathrm{Diff}(M)$ 下的协变性。考虑时空流形 $M$ 上的微分同胚 $\phi: M \to M$，张量场的变换为推前 (push-forward) 和拉回 (pull-back)：
\begin{equation}
    (\phi_* T)^{\mu_1\cdots\mu_k}_{\nu_1\cdots\nu_\ell}(p) 
    = \frac{\partial x'^{\mu_1}}{\partial x^{\rho_1}} \cdots 
      \frac{\partial x^{\nu_\ell}}{\partial x'^{\sigma_\ell}} \,
      T^{\rho_1\cdots\rho_k}_{\sigma_1\cdots\sigma_\ell}(\phi^{-1}(p)).
\end{equation}

\begin{note}[主动与被动观点]
广义协变性有两种等价的解读方式：
\begin{itemize}
    \item \textbf{被动微分同胚 (Passive diffeomorphism)}：物理系统不变，只是改变了描述它的坐标系。这相当于同一物理状态的不同坐标表示。
    \item \textbf{主动微分同胚 (Active diffeomorphism)}：坐标系不变，物理系统本身被"拖拽"到了流形上的不同位置。在广义相对论中，主动微分同胚是规范对称性——它将物理上不可区分的构型联系起来。
\end{itemize}
这两种观点的等价性是广义相对论规范结构的核心。
\end{note}

\subsection{Kretschmann 反对与物理解析}

1917年，Erich Kretschmann 提出了著名的反对意见：

\begin{quote}
    \textbf{Kretschmann 反对：} 任何物理理论——包括牛顿力学和狭义相对论——都可以通过引入适当的几何结构而被改写为"广义协变"的形式。因此广义协变性不是一个有物理内容的选择性原则。
\end{quote}

\begin{note}[Kretschmann 反对的解析]
Kretschmann 反对在技术上是正确的：通过足够精巧的数学操作，任何理论都可以写成广义协变的形式（例如，牛顿引力可以用 Newton--Cartan 几何表达）。然而，广义协变性在广义相对论中具有特殊的物理意义：
\begin{enumerate}
    \item \textbf{无背景结构 (No background structure)}：在广义相对论中，度规 $g_{\mu\nu}$ 本身是动力学变量，而牛顿力学和狭义相对论中存在先验的、非动力学的绝对时空结构（如 Minkowski 度规 $\eta_{\mu\nu}$）。
    \item \textbf{等效原理的实现}：广义协变性确保物理定律在局部惯性系中取狭义相对论的形式，这正是等效原理的数学表达。
    \item \textbf{规范对称性}：广义协变性是广义相对论的规范对称性，与之相关的约束方程（Hamiltonian 和 momentum 约束）是理论的动力学核心。
\end{enumerate}
因此，广义协变性的真正物理内容不在于"方程可以写成张量形式"，而在于\textbf{时空的几何结构本身是动力学的}——这是广义相对论与所有前相对论物理的根本区别。Anderson (1967) 和 Friedman (1983) 对这一问题有精辟的分析。
\end{note}

\section{从平直时空到弯曲时空}

\subsection{度规作为动力学场}

在狭义相对论中，时空是刚性、平坦的 Minkowski 时空，度规 $\eta_{\mu\nu} = \mathrm{diag}(-1,1,1,1)$ 是固定的、非动力学的背景。广义相对论的核心变革在于：\textbf{度规 $g_{\mu\nu}(x)$ 成为一个动力学场}——它既描述时空的几何，又描述引力。

\begin{definition}[度规的动力学角色]
在广义相对论中，度规张量 $g_{\mu\nu}(x)$ 扮演双重角色：
\begin{enumerate}
    \item \textbf{几何角色}：定义时空的度量结构——距离、角度、体积。
    \item \textbf{动力学角色}：作为引力场的势，其导数给出"引力场强"（Christoffel 符号），其二阶导数给出"潮汐力"（曲率张量）。
\end{enumerate}
\end{definition}

\begin{figure}[htbp]
\centering
\begin{tikzpicture}[scale=1.2]
    % 弯曲时空背景 (用波浪线暗示)
    \draw[gray!30, line width=2pt] plot[smooth, tension=0.8] coordinates {(-3.5,1.2) (-2.5,0.8) (-1.5,1.5) (-0.5,0.5) (0.5,1.3) (1.5,0.6) (2.5,1.4) (3.5,0.7)};
    \draw[gray!30, line width=2pt] plot[smooth, tension=0.8] coordinates {(-3.5,-0.2) (-2.5,0.3) (-1.5,-0.5) (-0.5,0.1) (0.5,-0.4) (1.5,0.2) (2.5,-0.3) (3.5,0.1)};
    
    % 局部惯性系（放大方框）
    \draw[blue, thick, fill=blue!5] (0.3, -0.2) rectangle (2.0, 1.1);
    \node[blue] at (1.15, 1.35) {局部惯性系 (LIF)};
    
    % 在局部惯性系中画Minkowski坐标网格
    \begin{scope}[shift={(0.5, 0.0)}]
        \draw[blue!60, very thin] (0.0,0.0) grid[step=0.3] (1.2,0.6);
        \draw[blue, ->] (0.6,0.3) -- (0.9,0.3) node[right] {$x$};
        \draw[blue, ->] (0.6,0.3) -- (0.6,0.6) node[above] {$t$};
    \end{scope}
    \node[blue, font=\small] at (1.8, -0.5) {$g_{\mu\nu} = \eta_{\mu\nu} + \mathcal{O}(x^2)$};

    % 自由下落粒子世界线
    \draw[red, thick, ->] (-3.0, -0.5) .. controls (-1.5, -0.8) and (-0.5, 0.8) .. (0.5, 1.0);
    \node[red] at (-1.8, 0.0) {自由下落粒子};
    
    % 全局弯曲时空标注
    \node at (-2.2, 1.6) {弯曲时空};
    \node at (-2.2, 1.2) {$(M, g_{\mu\nu})$};
    
    % 箭头指向局部惯性系
    \draw[->, thick, blue!70] (-1.5, 2.0) .. controls (0.0, 2.3) .. (1.15, 1.5);
    \node[blue] at (-0.5, 2.4) {局部等效于 Minkowski 时空};
    
\end{tikzpicture}
\caption{等效原理的几何表示。在弯曲时空 $(M, g_{\mu\nu})$ 的每一点，存在一个局部惯性系（Local Inertial Frame, LIF），在其中度规退化为 Minkowski 度规 $\eta_{\mu\nu}$ 到一阶精度。自由下落粒子在局部惯性系中沿直线运动，在全局弯曲时空中沿测地线运动。}
\label{fig:lif_curved}
\end{figure}

\subsection{自由运动从直线推广为测地线}

在狭义相对论中，自由粒子的运动方程为
\begin{equation}
    \frac{\dd^2 x^\mu}{\dd\tau^2} = 0,
\end{equation}
即粒子沿 Minkowski 时空中的直线运动。根据等效原理，在弯曲时空中，自由粒子应当沿"尽可能直"的路径运动——这就是\textbf{测地线}：

\begin{equation}
    \boxed{\frac{\dd^2 x^\mu}{\dd\tau^2} + \Gamma^\mu{}_{\nu\rho} \frac{\dd x^\nu}{\dd\tau} \frac{\dd x^\rho}{\dd\tau} = 0}.
\end{equation}

在局部惯性系中（$\Gamma^\mu{}_{\nu\rho} = 0$），测地线方程精确退化为 $\dd^2 x^\mu/\dd\tau^2 = 0$，满足等效原理的要求。

\section{最小耦合原理}

\subsection{原理的精确表述}

最小耦合原理 (Minimal Coupling Principle) 是将狭义相对论中的物理定律推广到弯曲时空的标准方法。

\begin{definition}[最小耦合原理 (Minimal Coupling Principle)]
将狭义相对论中成立的张量形式的物理定律推广到弯曲时空时，做以下替换：
\begin{enumerate}
    \item $\eta_{\mu\nu} \;\longrightarrow\; g_{\mu\nu}$ —— 将 Minkowski 度规替换为弯曲时空度规
    \item $\partial_\mu \;\longrightarrow\; \nabla_\mu$ —— 将普通偏导数替换为协变导数
    \item $\dd^4 x \;\longrightarrow\; \sqrt{-g} \, \dd^4 x$ —— 将不变体积元替换为广义协变的体积元
\end{enumerate}
除此之外，\textbf{不引入}任何包含曲率张量的额外耦合项。这就是"最小"的含义——只做维持广义协变性所必需的最小改动。
\end{definition}

\begin{note}[为什么是"最小"？]
"最小"意味着我们假定曲率与物质场的直接耦合不存在或可忽略。从有效场论的观点看，$R\phi^2$、$R_{\mu\nu}\partial^\mu\phi\partial^\nu\phi$ 等高阶项原则上可以存在，但它们被高能量标度（如 Planck 标度）所压低。在低能极限下，最小耦合是最自然的假设，且与所有现有实验一致。
\end{note}

\subsection{应用于 Klein-Gordon 标量场}

\textbf{平直时空：} Klein-Gordon 方程在 Minkowski 时空中为
\begin{equation}
    \eta^{\mu\nu}\partial_\mu\partial_\nu \phi - m^2\phi = 0.
\end{equation}
其作用量为
\begin{equation}
    S_{\text{KG}} = \int \dd^4 x \, \left[ -\frac{1}{2}\eta^{\mu\nu}\partial_\mu\phi\,\partial_\nu\phi - \frac{1}{2}m^2\phi^2 \right].
\end{equation}

\textbf{弯曲时空：} 应用最小耦合原理：
\begin{align}
    \eta^{\mu\nu} &\to g^{\mu\nu}, \\
    \partial_\mu &\to \nabla_\mu, \\
    \dd^4 x &\to \sqrt{-g}\,\dd^4 x.
\end{align}
注意对标量场 $\phi$，$\nabla_\mu\phi = \partial_\mu\phi$。作用量变为
\begin{equation}
    S_{\text{KG}} = \int \dd^4 x \, \sqrt{-g} \left[ -\frac{1}{2} g^{\mu\nu} \nabla_\mu\phi \,\nabla_\nu\phi - \frac{1}{2} m^2 \phi^2 \right].
\end{equation}

变分 $\delta S_{\text{KG}} / \delta\phi = 0$ 得到弯曲时空中的 Klein-Gordon 方程。利用
\begin{equation}
    \frac{1}{\sqrt{-g}} \partial_\mu\bigl(\sqrt{-g} g^{\mu\nu} \partial_\nu\phi\bigr) = g^{\mu\nu}\nabla_\mu\nabla_\nu\phi,
\end{equation}
有：
\begin{theorem}[弯曲时空 Klein-Gordon 方程]
\begin{equation}
    \boxed{g^{\mu\nu}\nabla_\mu\nabla_\nu \phi - m^2\phi = 0},
\end{equation}
或等价地
\begin{equation}
    \boxed{\frac{1}{\sqrt{-g}} \partial_\mu\bigl(\sqrt{-g}\, g^{\mu\nu} \partial_\nu\phi\bigr) - m^2\phi = 0}.
\end{equation}
协变 d'Alembert 算子定义为 $\Box_g \equiv g^{\mu\nu}\nabla_\mu\nabla_\nu$。
\end{theorem}

\subsection{应用于 Maxwell 电磁场}

\textbf{平直时空：} Maxwell 方程在 Minkowski 时空中的张量形式为
\begin{equation}
    \partial_\mu F^{\mu\nu} = J^\nu, \qquad 
    \partial_{[\mu}F_{\nu\rho]} = 0,
\end{equation}
其中 $F_{\mu\nu} = \partial_\mu A_\nu - \partial_\nu A_\mu$。

\textbf{弯曲时空：} 应用最小耦合原理。注意 $F_{\mu\nu}$ 是2-形式，其定义本身不依赖联络：
\begin{equation}
    F_{\mu\nu} = \nabla_\mu A_\nu - \nabla_\nu A_\mu = \partial_\mu A_\nu - \partial_\nu A_\mu,
\end{equation}
因为 Christoffel 符号的对称性 $\Gamma^\rho{}_{[\mu\nu]} = 0$ 使得协变导数与普通导数给出相同的结果。

Maxwell 方程推广为：
\begin{theorem}[弯曲时空 Maxwell 方程]
\begin{align}
    \nabla_\mu F^{\mu\nu} &= J^\nu, \label{eq:maxwell_cov} \\
    \nabla_{[\mu}F_{\nu\rho]} &= 0. \label{eq:maxwell_bianchi}
\end{align}
利用恒等式 $\nabla_\mu F^{\mu\nu} = \frac{1}{\sqrt{-g}}\partial_\mu(\sqrt{-g}F^{\mu\nu})$，(~\ref{eq:maxwell_cov})可写为
\begin{equation}
    \boxed{\frac{1}{\sqrt{-g}} \partial_\mu\bigl(\sqrt{-g}\,g^{\mu\alpha}g^{\nu\beta}F_{\alpha\beta}\bigr) = J^\nu}.
\end{equation}
\end{theorem}

\subsection{应用于完美流体}

\textbf{平直时空：} 完美流体的能量-动量张量为
\begin{equation}
    T^{\mu\nu} = (\rho + p) u^\mu u^\nu + p \, \eta^{\mu\nu},
\end{equation}
其中 $\rho$ 为能量密度，$p$ 为压强，$u^\mu$ 为四速度（满足 $u^\mu u_\mu = -1$）。守恒方程为 $\partial_\mu T^{\mu\nu} = 0$。

\textbf{弯曲时空：} 最小耦合原理给出：
\begin{theorem}[弯曲时空完美流体]
\begin{equation}
    \boxed{T^{\mu\nu} = (\rho + p) u^\mu u^\nu + p \, g^{\mu\nu}},
\end{equation}
其中 $u^\mu u_\mu = g_{\mu\nu}u^\mu u^\nu = -1$，守恒方程为
\begin{equation}
    \boxed{\nabla_\mu T^{\mu\nu} = 0}.
\end{equation}
展开协变导数给出广义相对论中的 Euler 方程（运动方程）和连续性方程。
\end{theorem}

\section{算例}

\subsection{例1：从等效原理计算引力红移}

\begin{example}[地球表面引力红移的定量预测]
\textbf{问题：} 利用等效原理计算地球表面高度差 $h=100\,\text{m}$ 处的引力红移。

\textbf{解：} 根据加速参考系的多普勒频移推导，频率变化为
\begin{equation}
    \frac{\Delta\nu}{\nu} = -\frac{gh}{c^2}.
\end{equation}
代入数值：
\begin{align}
    g &= 9.81 \,\text{m/s}^2, \\
    h &= 100 \,\text{m}, \\
    c &= 2.998 \times 10^8 \,\text{m/s}, \\[4pt]
    \frac{\Delta\nu}{\nu} &= -\frac{9.81 \times 100}{(2.998 \times 10^8)^2} 
    = -\frac{981}{8.988 \times 10^{16}} \approx -1.09 \times 10^{-14}.
\end{align}

\textbf{物理量纲检验：}
\begin{equation}
    [gh/c^2] = \frac{(\text{m/s}^2) \cdot \text{m}}{(\text{m/s})^2} = \text{无量纲}. \; \checkmark
\end{equation}

这一红移量相当于在 $10^{14}$ 个光子中少了一个光子的能量，其微小程度解释了为什么 Pound 和 Rebka 需要利用 M\"{o}ssbauer 效应的极高能量分辨率来进行测量。

\textbf{广义相对论验证：} 弱场度规为
\begin{equation}
    \dd s^2 = -\left(1 + \frac{2\Phi}{c^2}\right) c^2 \dd t^2 + \left(1 - \frac{2\Phi}{c^2}\right) \dd \vb*{x}^2.
\end{equation}
在地球表面，$\Phi = -GM_\oplus/r$。频率比由固有时比给出：
\begin{equation}
    \frac{\nu_{\text{obs}}}{\nu_{\text{em}}} = \sqrt{\frac{g_{00}(z_{\text{em}})}{g_{00}(z_{\text{obs}})}} 
    \approx 1 + \frac{\Phi_{\text{obs}} - \Phi_{\text{em}}}{c^2} = 1 - \frac{gh}{c^2}.
\end{equation}
与等效原理的推导完全一致。
\end{example}

\subsection{例2：标量场在弯曲时空中的波动方程}

\begin{example}[从最小耦合原理推导弯曲时空标量场方程]
\textbf{问题：} 从平直时空的无质量 Klein-Gordon 方程出发，应用最小耦合原理，推导弯曲时空中的波动方程，并验证在特定度规下的形式。

\textbf{解：} 
\textbf{步骤一——写出作用量：} 在 Minkowski 时空中，无质量标量场的作用量为
\begin{equation}
    S = -\frac{1}{2} \int \dd^4 x \, \eta^{\mu\nu} \partial_\mu \phi \, \partial_\nu \phi.
\end{equation}

\textbf{步骤二——最小耦合替换：}
\begin{equation}
    S = -\frac{1}{2} \int \dd^4 x \, \sqrt{-g} \, g^{\mu\nu} \partial_\mu \phi \, \partial_\nu \phi.
\end{equation}

\textbf{步骤三——变分：} $\delta S = 0$ 给出
\begin{align}
    \delta S &= -\frac{1}{2} \int \dd^4 x \, \sqrt{-g} \, \bigl[ g^{\mu\nu} \partial_\mu(\delta\phi) \partial_\nu\phi 
               + g^{\mu\nu} \partial_\mu\phi \, \partial_\nu(\delta\phi) \bigr] \\
             &= -\int \dd^4 x \, \sqrt{-g} \, g^{\mu\nu} \partial_\mu\phi \, \partial_\nu(\delta\phi) \\
             &= \int \dd^4 x \, \delta\phi \, \partial_\nu\bigl( \sqrt{-g} \, g^{\mu\nu} \partial_\mu\phi \bigr) = 0,
\end{align}
其中最后一步使用了分部积分并忽略边界项。因此
\begin{equation}
    \frac{1}{\sqrt{-g}} \partial_\nu\bigl(\sqrt{-g}\, g^{\mu\nu} \partial_\mu\phi\bigr) = 0.
\end{equation}

\textbf{步骤四——验证：} 设 FLRW 度规 $\dd s^2 = -\dd t^2 + a^2(t) \dd \vb*{x}^2$。此时
\begin{equation}
    \sqrt{-g} = a^3(t), \quad g^{00} = -1, \quad g^{ij} = a^{-2}(t) \delta^{ij}.
\end{equation}
代入方程：
\begin{align}
    \frac{1}{a^3} \partial_0\bigl( a^3 \cdot (-1) \cdot \partial_0\phi \bigr) 
    + \frac{1}{a^3} \partial_i\bigl( a^3 \cdot a^{-2} \delta^{ij} \partial_j\phi \bigr) &= 0, \\
    -\frac{1}{a^3} \partial_t(a^3 \dot\phi) + \frac{1}{a^2} \nabla^2 \phi &= 0, \\
    \ddot\phi + 3\frac{\dot a}{a} \dot\phi - \frac{1}{a^2} \nabla^2 \phi &= 0,
\end{align}
其中 $\dot\phi = \partial_t\phi$。这正是宇宙学中的标量场方程，$3H\dot\phi$ 项（$H=\dot a/a$ 为 Hubble 参数）代表宇宙膨胀的"摩擦"效应——这是最小耦合原理自动产生的结果。
\end{example}

\section*{本章小结}

本章从等效原理出发，逐步构建了广义相对论的物理基础。核心逻辑链如下：

\begin{center}
\begin{tikzpicture}[node distance=1.2cm, font=\small]
    \node[draw, rounded corners, fill=blue!10, align=center] (WEP) {$m_I=m_G$ \\ 弱等效原理};
    \node[draw, rounded corners, fill=red!10, right=of WEP, xshift=2.5cm, align=center] (EEP) {局部 SR 成立 \\ Einstein 等效原理};
    \node[draw, rounded corners, fill=green!10, right=of EEP, xshift=2.5cm, align=center] (SEP) {引力自能量也等效 \\ 强等效原理};
    
    \node[draw, rounded corners, fill=orange!10, below=of EEP, yshift=-0.8cm, align=center] (GC) {广义协变性 \\ 张量方程};
    \node[draw, rounded corners, fill=purple!10, below=of GC, yshift=-0.5cm, align=center] (MC) {最小耦合原理 \\ $\eta\to g,\;\partial\to\nabla$};
    \node[draw, rounded corners, fill=yellow!20, right=of MC, xshift=2.5cm, align=center] (GEOM) {弯曲时空物理定律};
    
    \draw[->, thick] (WEP) -- (EEP);
    \draw[->, thick] (EEP) -- (SEP);
    \draw[->, thick] (EEP) -- (GC);
    \draw[->, thick] (GC) -- (MC);
    \draw[->, thick] (MC) -- (GEOM);
    
    \node[below=0.3cm of WEP] {E\"{o}tv\"{o}s 实验};
    \node[below=0.3cm of EEP] {Pound--Rebka};
    \node[below=0.3cm of SEP] {LLR, 脉冲星};
\end{tikzpicture}
\end{center}

等效原理告诉我们：\textbf{引力在局部上与加速不可区分}，因此时空必须是弯曲的，而自由粒子沿测地线运动。广义协变性原理为这一几何描述提供了数学语言：物理定律必须以张量形式表达。最小耦合原理则提供了从平直时空物理定律构造弯曲时空物理定律的系统方法。三者共同构成了广义相对论的理论框架。在下一章中，我们将利用变分原理从 Einstein--Hilbert 作用量推导爱因斯坦场方程——它将物质分布与时空曲率定量地联系起来。
